\documentclass[11pt,a4paper]{article}
\usepackage[T1]{fontenc}
\usepackage{lmodern}
\usepackage[margin=27mm]{geometry}
\usepackage{amsmath,amssymb,amsthm,mathtools}
\usepackage{microtype,booktabs,array,longtable,enumitem}
\usepackage{xurl,hyperref,fancyhdr}
\hypersetup{hidelinks,pdftitle={A Unique Quartic Transition, Mobius Rigidity, and Two Radial Turning Points},pdfauthor={ProveIt research continuation}}
\setlength{\headheight}{14pt}
\pagestyle{fancy}\fancyhf{}
\fancyhead[L]{\small Real-order double polylogarithms}
\fancyhead[R]{\small ProveIt research continuation}
\fancyfoot[C]{\thepage}
\newtheorem{theorem}{Theorem}[section]
\newtheorem{lemma}[theorem]{Lemma}
\newtheorem{proposition}[theorem]{Proposition}
\newtheorem{corollary}[theorem]{Corollary}
\theoremstyle{definition}\newtheorem{definition}[theorem]{Definition}
\theoremstyle{remark}\newtheorem{remark}[theorem]{Remark}
\newcommand{\rbLi}{\operatorname{Li}}
\newcommand{\rbIm}{\operatorname{Im}}
\newcommand{\rbRe}{\operatorname{Re}}
\newcommand{\rbR}{\mathbb R}
\newcommand{\rbQ}{\mathbb Q}
\newcommand{\rbC}{\mathbb C}
\newcommand{\rbpa}{a_\dagger}
\newcommand{\rbpb}{b_\dagger}
\newcommand{\rbrr}{R_\dagger}
\newcommand{\rbmu}{\mu_\dagger}
\DeclareMathOperator{\diag}{diag}
\numberwithin{equation}{section}
\setlist{nosep}
\begin{document}
\begin{titlepage}
\centering
\vspace*{14mm}
{\large PROVEIT \quad / \quad ANALYSIS\par}
\vspace{13mm}
{\huge\bfseries A Unique Quartic Transition,\par}
\vspace{3mm}
{\huge\bfseries M\"obius Rigidity, and\par}
\vspace{3mm}
{\huge\bfseries Two Radial Turning Points\par}
\vspace{8mm}
{\Large for Real-Order Double Polylogarithms\par}
\vspace{14mm}
{\large Research continuation prepared for\par}
\vspace{2mm}
{\Large Vladimir Reshetnikov's ProveIt project\par}
\vspace{6mm}
{October 10, 2026\par}
\vfill
\begin{minipage}{0.91\textwidth}
\textbf{Abstract.}
For $F_{a,b}(z)=\sum_{n>m\ge1}z^n/(n^a m^b)$, write the upper angular
zero near the origin as $\cos\theta=\rho\eta_{a,b}(\rho)$ and expand
$\eta=\mu+K\rho^2+Q\rho^4+R\rho^6+\cdots$.
The audited ProveIt manuscript establishes a unique quadratic threshold
$a=A(b)$, but leaves the number of zeros of $q(b)=Q(A(b),b)$ and their
higher-order geometry open. We prove, with exact arithmetic certificates,
that $q$ has precisely one zero in $0<b<1$. Its location is enclosed to
40 decimal places, its derivative is positive, and the sixth-order radial
coefficient there is strictly negative. These facts yield a complete local
two-parameter bifurcation theorem: an open parameter region has two radial
turning points, a minimum followed by a maximum. We also give a fully
certified rational example, prove a fourth-root onset law, and classify
constant normalized-radius curves for $a\ge1$, $b>0$: only $(a,b)=(1,1)$
is possible. An all-order coefficient identity and its M\"obius transport
explain the rigidity and give an order-nine symmetry obstruction at the
transition. All analytic arguments are written out; the decisive finite
inequalities replay with the Python standard library alone.
\end{minipage}
\vfill
{\small Computer-assisted research manuscript. No proof-assistant formalization\par
or claim of literature-wide priority is made. Independent review is invited.\par}
\end{titlepage}
\tableofcontents
\newpage
\section{Research target and contribution}
\label{rbif:sec:overview}
This article continues the real-order and angular-zero part of
\emph{Polylogarithms and their Arithmetic Bridges}. The source audit is
pinned to the ProveIt repository state
\path{0707425e155707e53d2892e30c72e54d06c00e5f}.
The two directly relevant sections are
\path{05-real-local-radius.tex} and \path{05-real-turning.tex}
\cite{rbif:local,rbif:turning}. The former proves the sign of the first
radial coefficient and the existence of a unique fractional threshold.
The latter proves a maximum at inner order $1/2$, a minimum at inner order
$999/1000$, and the existence of an intervening quartic zero. Its closing
paragraph explicitly leaves the number of quartic zeros and the
higher-order radial behavior at those zeros open.

We resolve these finite-coefficient questions globally on $0<b<1$ and
then use the answer to obtain new analytic consequences for the actual
polylogarithm, rather than only for a Taylor polynomial. The main advances
are: a unique quartic transition; its nondegenerate sixth-order behavior;
a two-turning-point parameter region and a rational witness; an onset
exponent $1/4$ at the transition; and a constant-curve rigidity theorem.
The all-order identities in Section~\ref{rbif:sec:identities} are additional
exact functional and coefficient relations, not conjectured numerical
period reductions.

The proofs use standard analytic implicit-function arguments, elementary
one-variable bifurcation analysis, and explicit finite sums of real powers.
The global quartic sign theorem is computer-assisted. Its numerical inputs
are rational boxes whose validity is checked independently by integer
arithmetic. The box generator's floating-point root finder is not part of
the proof. The available literature on harmonic sums and polylogarithms
provides a broader algebraic setting \cite{rbif:abl}, while formal Lagrange
inversion supplies a standard coefficient-extraction tool
\cite{rbif:gessel}. We do not assert that a targeted literature search
establishes worldwide novelty.

\section{The analytic zero equation}
\label{rbif:sec:equation}
\subsection{Conventions and the local branch}
For $a\ge0$, $b\ge0$, define
\begin{equation}
 H_{n-1}^{(b)}=\sum_{m=1}^{n-1}m^{-b},\qquad
 p_n(a,b)=\frac{H_{n-1}^{(b)}}{n^a},\qquad
 F_{a,b}(z)=\sum_{n=2}^{\infty}p_n(a,b)z^n.
 \label{rbif:eq:def}
\end{equation}
For positive orders this is $\rbLi_{a,b}(z,1)$ in the convention
$n>m$. We work initially in $|z|<1$; no branch-cut convention is needed.
The value $b=0$ will only be used to close an interval in the
finite-coefficient argument. Since $0<p_n\le n-1$, the series and all
its $z$ derivatives converge normally on compact subdisks. Dependence on
$a,b$ is real analytic: on any compact complex neighborhood of a real
parameter pair, the coefficients have a uniform polynomial growth bound.

Let $z=\rho e^{i\theta}$, with $0<\rho<1$ and $0<\theta<\pi$, and set
\begin{equation}
 t=\rho^2,\qquad y=\frac{\cos\theta}{\rho}.
 \label{rbif:eq:ty}
\end{equation}
The manuscript proves a unique upper-arc angular zero throughout its
stated real-order domain \cite{rbif:cartesian}. Our main local arguments
need only the branch constructed below; the explicit witness will also
have its own certified existence strip.

\begin{theorem}[Exact angular generating equation]
\label{rbif:thm:generating}
Put
\begin{equation}
 P_k(y)=\sum_{j=0}^{k+1}(-1)^j\binom{k+1}{j}
             (2y)^{k+1-j}p_{k+2+j}\qquad(k\ge0).
 \label{rbif:eq:Pk}
\end{equation}
Then, near $t=0$,
\begin{equation}
 \Psi(t,y):=\sum_{k\ge0}P_k(y)t^k
 =\frac{\rbIm F_{a,b}(\rho e^{i\theta})}{\rho^3\sin\theta}
 \quad\text{when }t=\rho^2,\quad\cos\theta=\rho y.
 \label{rbif:eq:Psi}
\end{equation}
There is a unique real-analytic solution $y=\Phi(a,b,t)$ near $t=0$,
with
\begin{equation}
 \Phi(a,b,0)=\mu(a,b):=\frac{p_3}{2p_2}.
 \label{rbif:eq:mu}
\end{equation}
Consequently $\eta_{a,b}(\rho)=\Phi(a,b,\rho^2)$ is even and real
analytic near zero.
\end{theorem}
\begin{proof}
The identity
$\sin(n\theta)/\sin\theta=U_{n-1}(\cos\theta)$ and the elementary
expansion
\[
 U_{n-1}(x)=\sum_{j=0}^{\lfloor(n-1)/2\rfloor}
          (-1)^j\binom{n-1-j}{j}(2x)^{n-1-2j}
\]
give the right side of \eqref{rbif:eq:Psi}. A term with indices $n,j$
has radial power $\rho^{2n-4-2j}=t^{n-2-j}$. Regrouping at
$k=n-2-j$ gives \eqref{rbif:eq:Pk}. The only apparent negative power,
from $n=2$, cancels because $U_1(\rho y)=2\rho y$.
For $|y|\le M$,
\begin{equation}
 |P_k(y)|\le (2k+3)(1+2M)^{k+1}.
 \label{rbif:eq:majorant}
\end{equation}
Thus regrouping is justified when $|t|(1+2M)<1$. At $t=0$ the
equation is $2p_2y-p_3=0$, and its $y$ derivative is $2p_2>0$.
The analytic implicit-function theorem proves all remaining claims.
\end{proof}

The same equation can be derived without trigonometry. If $z,w$ are
the roots of $X^2-2tyX+t=0$, then
\[
 \Psi(t,y)=\frac{F(z)-F(w)}{t(z-w)}.
\]
This is an identity of symmetric power series, with the removable factors
understood at the origin. For the actual angular arc, $w=\overline z$.
This formulation makes the algebraic symmetry in later sections natural.

\subsection{A usable tail bound}
Let $L=1+2M$, $q=L|t|<1$, and truncate after $P_Nt^N$. Summing
\eqref{rbif:eq:majorant} gives
\begin{equation}
 \left|\Psi(t,y)-\sum_{k=0}^NP_k(y)t^k\right|
 \le Lq^{N+1}\left(\frac{2N+5}{1-q}
                         +\frac{2q}{(1-q)^2}\right).
 \label{rbif:eq:tail}
\end{equation}
This bound does not exploit cancellation and therefore remains valid for
all nonnegative orders. With $n=N+1$, a bound for the $t$ derivative of
the tail is
\begin{equation}
 L^2q^{n-1}\left\{
 \frac{2n^2+3n}{1-q}+\frac{(4n+3)q}{(1-q)^2}
                  +\frac{2q(1+q)}{(1-q)^3}\right\}.
 \label{rbif:eq:tailt}
\end{equation}
Both inequalities follow by summing a geometric series and its first
two derivatives. For real $1/4\le y\le3/4$, the complex $y$ disk of
radius $1/4$ lies in $|y|\le1$. Cauchy's estimate therefore bounds the
$y$ derivative of the tail by four times \eqref{rbif:eq:tail} with $M=1$.
These explicit estimates are used in the rational witness, so its
conclusions concern the infinite function.

\section{All-order identities and M\"obius transport}
\label{rbif:sec:identities}
\subsection{A coefficient-extraction identity}
Write
\begin{equation}
 \Phi(a,b,t)=\mu+\sum_{N\ge1}c_Nt^N.
 \label{rbif:eq:cN}
\end{equation}
The first terms $c_1,c_2,c_3$ will be denoted $K,Q,R$.

\begin{proposition}[All-order reversion formula]
\label{rbif:prop:reversion}
For every $N\ge1$,
\begin{equation}
 c_N=\sum_{m=1}^N\frac{(-1)^m}{m(2p_2)^m}
 [u^{m-1}]\!
 \sum_{\substack{k_1+\cdots+k_m=N\\k_1,\ldots,k_m\ge1}}
       \prod_{r=1}^m P_{k_r}(\mu+u).
 \label{rbif:eq:reversion}
\end{equation}
In particular, $c_N$ depends only on $p_2,\ldots,p_{2N+3}$.
\end{proposition}
\begin{proof}
Set $v=\Phi-\mu$. The equation is
$v=-(2p_2)^{-1}\sum_{k\ge1}P_k(\mu+v)t^k$.
Introduce an auxiliary multiplier $s$ on the right side and apply the
formal identity
$[s^m]v=m^{-1}[u^{m-1}]H(t,u)^m$ to $v=sH(t,v)$.
For completeness, this identity follows by substituting
$s=u/H(t,u)$ in formal coefficient extraction, or by differentiating
powers and comparing their residues; it is the standard Lagrange
inversion formula \cite{rbif:gessel}. Each factor of $H$ is divisible
by $t$, so after extracting $t^N$ only $m\le N$ survive. Setting $s=1$
is thus legitimate in the $t$-adic topology, and gives
\eqref{rbif:eq:reversion}. Formula \eqref{rbif:eq:Pk} gives the stated
finite dependence. Analytic convergence follows independently from
Theorem~\ref{rbif:thm:generating}.
\end{proof}

\subsection{The sixth-order radial coefficient}
Let $A(y)=P_1(y)$, $B(y)=P_2(y)$, $C(y)=P_3(y)$. Explicitly,
\begin{align}
 A(y)&=4p_3y^2-4p_4y+p_5,\label{rbif:eq:ABC1}\\
 B(y)&=8p_4y^3-12p_5y^2+6p_6y-p_7,\\
 C(y)&=16p_5y^4-32p_6y^3+24p_7y^2-8p_8y+p_9.
 \label{rbif:eq:ABC3}
\end{align}
Substitution into $\Psi(t,\Phi(t))=0$ gives the exact identities
\begin{align}
 K&=-\frac{A(\mu)}{2p_2},\label{rbif:eq:K}\\
 Q&=-\frac{A'(\mu)K+B(\mu)}{2p_2},\label{rbif:eq:Q}\\
 R&=-\frac{A'(\mu)Q+4p_3K^2+B'(\mu)K+C(\mu)}{2p_2}.
 \label{rbif:eq:R}
\end{align}
The formulas for $K,Q$ agree with the audited manuscript. The coefficient
$R$ is the next obstruction needed at its unresolved transition.
At a simultaneous zero of $K,Q$ it simplifies to
\begin{equation}
 R=-\frac{C(\mu)}{2p_2}.
 \label{rbif:eq:Rdeg}
\end{equation}
These expressions are exact in arbitrary real orders; integer relations
or numerical fitting are not involved.

For efficient verification let $r_n=p_n/p_2=(2/n)^aH_{n-1}^{(b)}$ and
$s=r_3$. Then $\mu=s/2$ and
\begin{align}
 K&=\frac{2sr_4-s^3-r_5}{2},\label{rbif:eq:ratioK}\\
 Q&=-2(s^2-r_4)K
       -\frac{r_4s^3-3r_5s^2+3r_6s-r_7}{2}.\label{rbif:eq:ratioQ}
\end{align}
The corresponding ratio formula for $R$ is
\begin{align}
 R={}&-2(s^2-r_4)Q-2sK^2
       -3(r_4s^2-2r_5s+r_6)K\notag\\
 &-\frac{r_5s^4-4r_6s^3+6r_7s^2-4r_8s+r_9}{2}.
 \label{rbif:eq:ratioR}
\end{align}
Avoiding repeated interval differentiation of $p_2^{-1}$ substantially
reduces enclosure widths.

\subsection{An exact M\"obius identity}
For a real constant $c$, define the involution and the local coordinate
\begin{equation}
 T_c(z)=-\frac{z}{1-2cz},\qquad
 w=\frac{z}{1-cz},\qquad z=\frac{w}{1+cw}.
 \label{rbif:eq:mobius}
\end{equation}
The involution sends $w$ to $-w$. Write
$G_c(w)=F(w/(1+cw))=\sum_{n\ge2}d_n(c)w^n$.

\begin{theorem}[Transport of the angular coefficients]
\label{rbif:thm:transport}
Let $F(z)=\sum_{n\ge2}p_nz^n$ be holomorphic near zero, with $P_k$ defined by
\eqref{rbif:eq:Pk}, the following are identities of convergent local
power series:
\begin{align}
 d_{2j+3}(c)
   &=\sum_{k=0}^j(-1)^{k+1}\binom{j+1}{k+1}
                      c^{2(j-k)}P_k(c),\label{rbif:eq:transport}\\
 G_c(w)-G_c(-w)
   &=-\frac{2w^3}{(1-c^2w^2)^2}
      \Psi\!\left(-\frac{w^2}{1-c^2w^2},c\right).
 \label{rbif:eq:mobiusPsi}
\end{align}
\end{theorem}
\begin{proof}
Take $w=iv$ for small real $v$. Then
$z=(cv^2+iv)/(1+c^2v^2)$, so
$t=|z|^2=v^2/(1+c^2v^2)$ and $\rbRe z=c|z|^2$.
Equation \eqref{rbif:eq:Psi} yields
\[
 \rbIm G_c(iv)=\frac{v^3}{(1+c^2v^2)^2}
        \Psi\!\left(\frac{v^2}{1+c^2v^2},c\right).
\]
For real coefficients, twice this imaginary part equals
$(G_c(iv)-G_c(-iv))/i$. Replacing $iv$ by $w$ gives
\eqref{rbif:eq:mobiusPsi}, first on an accumulating set and then as a
holomorphic identity. Both sides are polynomial-linear in each $p_n$
at any fixed degree, so the same identity holds for complex coefficients.
Expanding $(1+c^2v^2)^{-k-2}$ and comparing coefficients gives
\eqref{rbif:eq:transport}. Alternatively the latter follows directly
from
$d_n(c)=\sum_{m=2}^n\binom{n-1}{m-1}(-c)^{n-m}p_m$.
\end{proof}

\begin{corollary}[Structure of constant normalized-radius curves]
\label{rbif:cor:flatstructure}
Suppose $F$ has real coefficients, $p_2>0$, and its local angular branch
has $\eta(\rho)\equiv c$. This happens if and only if
\begin{equation}
 F(T_c(z))=F(z),
 \quad\text{equivalently}\quad
 F(z)=\mathcal H\!\left(\frac{z^2}{(1-cz)^2}\right)
 \label{rbif:eq:flatstructure}
\end{equation}
for a holomorphic germ $\mathcal H$ with real coefficients and
$\mathcal H(0)=0$. Necessarily $c=p_3/(2p_2)$.
If $\mathcal H(u)=\sum_{j\ge1}h_ju^j$, its coefficient identities are
\begin{equation}
 p_n=\sum_{j=1}^{\lfloor n/2\rfloor}
           h_j\binom{n-1}{2j-1}c^{n-2j}\qquad(n\ge2).
 \label{rbif:eq:flatcoeff}
\end{equation}
\end{corollary}
\begin{proof}
Constancy means $\Psi(t,c)=0$ identically. The invertible substitution in
\eqref{rbif:eq:mobiusPsi} makes this equivalent to $G_c$ being even.
An even holomorphic germ has the form $\mathcal H(w^2)$, and the converse
is immediate. The coefficient of $z^3$ determines $c$, and expansion of
$(1-cz)^{-2j}$ proves \eqref{rbif:eq:flatcoeff}.
\end{proof}

\section{The quadratic threshold and its scope}
\label{rbif:sec:quadratic}
We recall the manuscript's quadratic-threshold result, including the
inequalities needed later for rigidity. A compact proof is included to
make the dependency explicit; this result is credited to the audited
manuscript, not presented as a new contribution \cite{rbif:local}.

\begin{theorem}[Established quadratic classification]
\label{rbif:thm:quadratic}
For $a\ge1$, $b>0$:
\begin{enumerate}[label=(\roman*)]
\item $K(a,b)<0$ for every $a\ge2$;
\item $K(1,b)$ has the sign of $1-b$;
\item if $b\ge1$ and $a>1$, then $K(a,b)<0$;
\item for each $0<b<1$ there is exactly one $A(b)\in(1,2)$ with
$K(A(b),b)=0$. Below this threshold $K>0$, and above it $K<0$.
\end{enumerate}
The threshold extends analytically to a neighborhood of each endpoint
of $[0,1]$, with $A(1)=1$, and $K_a(A(b),b)<0$.
\end{theorem}
\begin{proof}
Put $x=2^{-b}$, $y=3^{-b}$, $U=1+x$, $V=1+x+y$,
$W=1+x+y+x^2$. Direct substitution in \eqref{rbif:eq:K} gives
\begin{equation}
 K=-\tfrac12(2/5)^a G_a,\qquad
 G_a=W+U^3(20/27)^a-2UV(5/6)^a.
 \label{rbif:eq:Ga}
\end{equation}
Here the subscript in $G_a$ is an order parameter, not a derivative.
Since $x^2\le y$ and $0\le x\le1$,
$U^2/V\le4/3$. The rational logarithm bounds
\[
 \log(27/20)<\frac{7273}{24000},\qquad
 \log(6/5)>\frac9{50}
\]
therefore imply
\[
 \frac{U^2}{2V}\frac{\log(27/20)}{\log(6/5)}
 <\frac{7273}{6480}<\frac98\le(9/8)^a.
\]
This is exactly the inequality $\partial_aG_a>0$ for $a\ge1$.
The displayed logarithm estimates follow from the alternating expansion
of $\log(1+u)$ with its next-term bound.

For $a=2$, use $y\le x^{3/2}$ and set $x=t^2$. The coefficient of $y$
in $G_2$ is negative, and substitution gives the lower bound
\[
 G_2\ge\frac{800t^6-2025t^5+1833t^4-567t^3-192t^2+233}{1458}.
\]
Its degree-six Bernstein coefficients on $[0,1]$ are
\[
 \left(\frac{233}{1458},\frac{233}{1458},\frac{367}{2430},
 \frac{665}{5832},\frac{55}{486},\frac{95}{1458},\frac{41}{729}\right).
\]
They are all positive. Thus $G_2>0$, and monotonicity gives assertion (i).

At $a=1$,
\begin{equation}
 27G_1=20x^3+42x^2-3x+2-(45x+18)y.
 \label{rbif:eq:G1}
\end{equation}
For $b>1$, set $t=\sqrt{2x}<1$. The bound $y\le t^3/3$ gives
$G_1\ge(1-t)P_5(t)$, where
\[
 P_5(t)=\frac{-5t^5+10t^4-11t^3+t^2+4t+4}{54}.
\]
Its Bernstein coefficients are
$2/27,4/45,19/180,14/135,1/10,1/18$, so $G_1>0$.
For $0<b<1$, write $p=\log_2 3$. Since $11/7<p<2$, Taylor's
integral formula at $x=1/2$ gives
\[
 x^p\ge\frac13+\frac{22}{21}(x-1/2)+\frac{22}{49}(x-1/2)^2.
\]
Indeed $p(p-1)x^{p-2}>44/49$ on $1/2<x<1$.
Using the negative coefficient of $y$ in \eqref{rbif:eq:G1} yields
\[
 G_1\le-\frac{(2x-1)(10x^2-337x+334)}{2646}<0.
\]
The quadratic in parentheses decreases on $[1/2,1]$ and equals $7$ at
one. Direct substitution gives $G_1=0$ at $b=1$.
These facts prove (ii)--(iv), using $\partial_aG_a>0$.
At a threshold, \eqref{rbif:eq:Ga} gives $K_a<0$.
At $b=0$, $G_1=-2/27<0$ and $G_2>0$, so the same root persists;
at $b=1$ the root is $a=1$. The implicit-function theorem proves the
claimed endpoint extensions.
\end{proof}

\section{Exactly one quartic transition}
\label{rbif:sec:quartic}
Define on the closed interval $[0,1]$
\begin{equation}
 q(b)=Q(A(b),b).
 \label{rbif:eq:qdef}
\end{equation}
The endpoint extensions in Theorem~\ref{rbif:thm:quadratic} make this a
real-analytic function near every point of the interval.

\begin{theorem}[Global quartic sign classification]
\label{rbif:thm:quartic}
There is exactly one $\rbpb\in(0,1)$ with $q(\rbpb)=0$.
Moreover,
\begin{equation}
 q(b)<0\quad(0<b<\rbpb),\qquad
 q(b)>0\quad(\rbpb<b<1),\qquad q'(\rbpb)>0.
 \label{rbif:eq:qsigns}
\end{equation}
The zero and its corresponding outer order $\rbpa=A(\rbpb)$ satisfy
\begin{equation}
 \begin{gathered}
 0.9974937898734204210746055930669544902305<\rbpb,\\
 \rbpb<0.9974937898734204210746055930669544902306,
 \end{gathered}\label{rbif:eq:bbox}
\end{equation}
\begin{equation}
 \begin{gathered}
 1.0014301809614951381293517293887015374543<\rbpa,\\
 \rbpa<1.0014301809614951381293517293887015374544.
 \end{gathered}\label{rbif:eq:abox}
\end{equation}
The inequalities are certified enclosures, not rounded root estimates.
\end{theorem}

The finite certificate underlying the global part of the proof is the
following stronger pair of derivative inequalities.
\begin{lemma}[Certified derivative separation]
\label{rbif:lem:separation}
The exact replay proves
\begin{equation}
 q'(b)>10^{-6}\quad(0\le b\le9/10),\qquad
 q''(b)<-\frac{29}{10^6}\quad(9/10\le b\le1).
 \label{rbif:eq:separation}
\end{equation}
It also proves $q(9/10)<0<q(999/1000)$, while $q(1)=0$ exactly.
\end{lemma}
\begin{proof}[Proof and finite verification specification]
The computational claim is reduced to finitely many explicit rational
inequalities as follows. Partition $[0,9/10]$ into $2048$ equal cells
and $[9/10,1]$ into $256$ equal cells. At every endpoint and midpoint,
provide a rational interval $[a_-,a_+]$ for $A(b)$ and verify
$K(a_-,b)>0>K(a_+,b)$. There are $4609$ distinct such intervals.
For a cell $[u,v]$ form the rectangle
\[
 \mathcal B=[a_-(v),a_+(u)]\times[u,v].
\]
The replay verifies $K_a<0$ and $K_b<0$ throughout this rectangle.
The signs at the two opposite corners, together with $K_b<0$, are
therefore barriers for the entire graph of $A$ on the cell. No unverified
interpolation of the graph is used.

All partial derivatives through degree three are computed by automatic
differentiation of finite exponential sums, using exact outward dyadic
intervals. It is convenient to replace $Q$ by
\begin{equation}
 Q_0=-\tfrac12(r_4s^3-3r_5s^2+3r_6s-r_7),\qquad s=r_3.
 \label{rbif:eq:Q0}
\end{equation}
Since $Q=Q_0-2(s^2-r_4)K$, the restrictions of $Q$ and $Q_0$ to the
threshold, and all their total $b$ derivatives there, agree exactly.
Implicit differentiation gives
\begin{align}
 A'&=-K_b/K_a,\notag\\
 A''&=-(K_{bb}+2K_{ab}A'+K_{aa}(A')^2)/K_a,\label{rbif:eq:implicitdiff}\\
 q'&=(Q_0)_b+(Q_0)_aA',\notag\\
 q''&=(Q_0)_{bb}+2(Q_0)_{ab}A'+(Q_0)_{aa}(A')^2+(Q_0)_aA''.
 \notag
\end{align}
A further differentiation gives $q'''$, with its explicit formula in
Appendix~\ref{rbif:app:verification}.

Let $m=(u+v)/2$ and $h=(v-u)/2$. On a first-part cell use the rigorous
centered bound
\[
 q'(b)\ge \inf q'(m)-h\sup_{[u,v]}|q''|.
\]
On a second-part cell use
\[
 q''(b)\le \sup q''(m)+h\sup_{[u,v]}|q'''|.
\]
The first lower bounds all exceed $10^{-6}$; the negatives of the second
upper bounds all exceed $29/10^6$. The exact minimum margins are recorded
as integer numerators over $2^{256}$ in \path{global_verified.json}.
Their decimal diagnostics are $1.1282650208766\cdot10^{-6}$ and
$2.9857032148394\cdot10^{-5}$, respectively. Only the integer comparisons
are used in the proof.

For the endpoint signs the same arithmetic gives
$q(9/10)\approx-1.86117461963559\cdot10^{-7}$ and
$2.92779196\cdot10^{-11}<q(999/1000)<2.92779197\cdot10^{-11}$.
At $(a,b)=(1,1)$, direct rational substitution in the finite coefficient
formulas gives $K=Q=0$, hence $q(1)=0$.
The programs \path{verify_global.py} and \path{verify_local.py} check
all the stated inequalities. Elementary-function tails, arithmetic rules,
input boxes, full-cover checks, and executable replay instructions are
included in the package and described in
Appendix~\ref{rbif:app:verification}. This is a finite exact proof
certificate, not a sampling argument.
\end{proof}

\begin{proof}[Proof of Theorem~\ref{rbif:thm:quartic}]
By Lemma~\ref{rbif:lem:separation}, $q$ is increasing on $[0,9/10]$
and is negative at its right endpoint. Thus it has no zero there.
On $[9/10,1]$ it is strictly concave. The opposite signs at $9/10$
and $999/1000$ give an interior zero, and $q(1)=0$ gives a second zero
at the endpoint. A strictly concave function cannot have three distinct
zeros: at the middle one it would lie strictly above the chord between
the outer two. Thus there is exactly one interior zero. Strict concavity
also makes $q$ positive between that zero and one, and its derivative at
the interior zero is positive. The preceding signs determine the negative
side. The exact localization is proved independently in the next section;
its unique local zero must be this globally unique zero.
\end{proof}

\begin{corollary}[Complete quartic signs at the quadratic threshold]
\label{rbif:cor:localfamily}
For every fixed $0<b<\rbpb$, decreasing $a$ slightly below $A(b)$
creates a unique small-radius strict maximum. For every fixed
$\rbpb<b<1$, increasing $a$ slightly above $A(b)$ creates a unique
small-radius strict minimum. The assertions concern a sufficiently small
radius interval, allowed to depend on $b$.
\end{corollary}
\begin{proof}
With $U=\partial_t\Phi$, at $(A(b),b,0)$ one has
$U=0$, $U_t=2q(b)\ne0$, and $U_a=K_a<0$.
The analytic implicit-function theorem produces one zero of $U$ in $t$.
Its side, derivative sign, and extremum type follow from the displayed
signs. Since $\partial_\rho\eta=2\rho U$, the same sign changes hold
in the physical radius. This generalizes the two separately certified
inner orders in the earlier manuscript.
\end{proof}

\section{The degenerate point is nondegenerate at the next order}
\label{rbif:sec:point}
Set
\[
 \rbrr=R(\rbpa,\rbpb),\qquad \rbmu=\mu(\rbpa,\rbpb),\qquad
 J_\dagger=\det\frac{\partial(K,Q)}{\partial(a,b)}(\rbpa,\rbpb).
\]
\begin{theorem}[Exact nondegeneracy certificate]
\label{rbif:thm:nondeg}
At the point in \eqref{rbif:eq:bbox}--\eqref{rbif:eq:abox},
\begin{align}
 -3.368744536\cdot10^{-11}&<\rbrr<-3.368744535\cdot10^{-11},
 \label{rbif:eq:Rbox}\\
 -8.315139844\cdot10^{-10}&<J_\dagger<-8.315139843\cdot10^{-10},
 \label{rbif:eq:Jbox}\\
 4.870284716\cdot10^{-8}&<q'(\rbpb)<4.870284717\cdot10^{-8},
 \label{rbif:eq:qpbox}\\
 -0.017073210967&<K_a(\rbpa,\rbpb)<-0.017073210965.
 \label{rbif:eq:Kabox}
\end{align}
In particular, $(a,b)\mapsto(K,Q)$ is a real-analytic local coordinate
system at this point. The radial branch satisfies
\begin{equation}
 \eta_{\rbpa,\rbpb}(\rho)=\rbmu+\rbrr\rho^6+O(\rho^8),
 \label{rbif:eq:degexp}
\end{equation}
and is strictly decreasing for sufficiently small positive $\rho$.
\end{theorem}
\begin{proof}
Here is the existence certificate, not just a residual estimate.
Let $X$ be the rational rectangle in
\eqref{rbif:eq:bbox}--\eqref{rbif:eq:abox}, in $(a,b)$ order, and
let $c$ be its exact rational midpoint. Put $f=(K,Q)$.
Evaluate $Df(c)$ by intervals and take the exact rational inverse $C$
of its entrywise midpoint. Interval evaluation on the whole rectangle
then gives
\[
 H=I-C\,Df(X),\qquad
 f\text{-iteration image }\mathcal N
 =c-Cf(c)+H(X-c).
\]
The verifier proves $\mathcal N\subset\operatorname{int}X$ and
\begin{equation}
 \|H\|_\infty<3\cdot10^{-32}<1.
 \label{rbif:eq:contraction}
\end{equation}
Indeed its diagnostic upper bound is
$2.6668469993303\cdot10^{-32}$; the actual comparison is rational.
For $T(x)=x-Cf(x)$, the integral mean-value formula shows that
$T(X)\subseteq\mathcal N$. The same derivative enclosure makes $T$
a contraction on $X$. Iteration is a Cauchy sequence in the complete
closed rectangle and converges to a unique fixed point. Since $C$ is
invertible, this fixed point has $K=Q=0$.

The exact inverse matrix, image intervals, and derivative norm are in
\path{local_verified.json}. Direct interval substitution into
\eqref{rbif:eq:ratioR} and the derivatives of
\eqref{rbif:eq:ratioK}--\eqref{rbif:eq:ratioQ} proves
\eqref{rbif:eq:Rbox}--\eqref{rbif:eq:Kabox}. The identity
\[
 q'(b)=\frac{K_aQ_b-K_bQ_a}{K_a}\bigg|_{(A(b),b)}
\]
provides an additional consistent route to the sign in
\eqref{rbif:eq:qpbox}. Analytic inversion follows from
$J_\dagger\ne0$. Finally, \eqref{rbif:eq:degexp} follows from
$K=Q=0$ and \eqref{rbif:eq:R}; its derivative is
$6\rbrr\rho^5+O(\rho^7)<0$ for small positive radius.
\end{proof}

\begin{remark}[Why a high-precision residual was not enough]
The determinant has magnitude about $8.3\cdot10^{-10}$, so the two
coefficient equations are nearly tangent in the original parameters.
The sixth-order coefficient is smaller still. A small residual alone
would not establish either existence or the required signs. The
contraction and full-rectangle derivative bounds address both issues.
\end{remark}

\section{The full local two-parameter bifurcation}
\label{rbif:sec:bifurcation}
By Theorem~\ref{rbif:thm:nondeg}, use the exact coefficient coordinates
\begin{equation}
 \kappa=K(a,b),\qquad \lambda=Q(a,b)
 \label{rbif:eq:coords}
\end{equation}
near $(\rbpa,\rbpb)$. Write $\Phi(\kappa,\lambda,t)$ for the analytic
radial branch in these coordinates and put $U=\partial_t\Phi$.
The coefficient definitions give
\begin{equation}
 U(\kappa,\lambda,t)=\kappa+2\lambda t
        +3R(\kappa,\lambda)t^2+O(t^3),\qquad
 R(0,0)=\rbrr<0.
 \label{rbif:eq:U}
\end{equation}
In particular $U(\kappa,\lambda,0)=\kappa$ and
$U_t(\kappa,\lambda,0)=2\lambda$ exactly.

\begin{theorem}[Complete local turning-point classification]
\label{rbif:thm:unfolding}
There exist $t_0>0$, a parameter neighborhood $\mathcal V$ of $(0,0)$,
and a real-analytic function $\varphi$ near zero such that
\begin{equation}
 \varphi(0)=\varphi'(0)=0,\qquad
 \varphi(\lambda)=\frac{\lambda^2}{3\rbrr}+O(\lambda^3).
 \label{rbif:eq:fold}
\end{equation}
Within $\mathcal V$, the following list gives every critical point of
$\eta(\rho)$ in $0<\rho<\sqrt{t_0}$.
\begin{center}
\begin{tabular}{@{}p{0.45\textwidth}p{0.46\textwidth}@{}}
\toprule
Parameter condition & Critical-point behavior\\
\midrule
$\lambda\le0$, $\kappa\le0$ & No critical point.\\[1mm]
$\kappa>0$ (either sign of $\lambda$) & Exactly one, a strict maximum.\\[1mm]
$\lambda>0$, $\kappa<\varphi(\lambda)$ & No critical point.\\[1mm]
$\lambda>0$, $\kappa=\varphi(\lambda)$ & One stationary point, not an extremum.\\[1mm]
$\lambda>0$, $\varphi(\lambda)<\kappa<0$
 & Exactly two: a strict minimum, then a strict maximum.\\[1mm]
$\lambda>0$, $\kappa=0$ & Exactly one positive-radius critical point, a strict maximum.\\
\bottomrule
\end{tabular}
\end{center}
On the fold $\kappa=\varphi(\lambda)$, the stationary point has
$\eta'=\eta''=0$ and $\eta'''<0$. The two-critical-point region is open
and nonempty in the original $(a,b)$ plane.
\end{theorem}
\begin{proof}
Since $U_{tt}(0,0,0)=6\rbrr<0$, choose a sufficiently small closed
parameter and $t$ neighborhood on which $U_{tt}<0$ and $U_\kappa>0$.
Choose $t_0>0$ inside it. At the base parameters,
$U(t_0)<0$ and $U_t(t_0)<0$; by shrinking the parameter neighborhood
these inequalities hold uniformly. Thus $U$ is strictly concave as a
function of $t$, with a negative value at the upper boundary.

If $\lambda\le0$, then $U_t(0)=2\lambda\le0$, so $U$ is strictly
decreasing for positive $t$. Its initial value $\kappa$ immediately
gives the first two cases.
If $\lambda>0$, then $U_t(0)>0>U_t(t_0)$, so strict concavity gives
one interior maximum of $U$, at a point $T(\kappa,\lambda)$.
The implicit-function theorem applied to $U_t=0$ constructs $T$
analytically near the origin, even if its analytic extension is negative
when $\lambda<0$. Differentiation at the origin gives
\[
 T_\kappa(0,0)=0,\qquad T_\lambda(0,0)=-\frac1{3\rbrr}.
\]
Set $H(\kappa,\lambda)=U(\kappa,\lambda,T(\kappa,\lambda))$.
Then $H_\kappa(0,0)=1$, so $H=0$ is an analytic curve
$\kappa=\varphi(\lambda)$. Along this curve,
\[
 T=-\frac{\lambda}{3\rbrr}+O(\lambda^2),\qquad
 0=H=\varphi(\lambda)-\frac{\lambda^2}{3\rbrr}+O(\lambda^3).
\]
This proves \eqref{rbif:eq:fold}; in particular $\varphi(\lambda)<0$
for small nonzero $\lambda$.
Since $H_\kappa>0$ in a smaller neighborhood, the sign of the maximum
of $U$ is the sign of $\kappa-\varphi(\lambda)$. Together with
$U(0)=\kappa$ and $U(t_0)<0$, this gives every remaining row of the table.
Strict concavity precludes any additional root in this interval.

At a simple root, $\eta_{\rho\rho}=4tU_t$. Thus the upward crossing
is a strict minimum and the downward crossing is a strict maximum.
At a double root of $U$, one has $U=U_t=0$ and
\[
 \eta_{\rho\rho\rho}=8\rho^3U_{tt}<0.
\]
There is no change of sign in $\eta'$ there, hence no local extremum.
Finally, analytic inversion of \eqref{rbif:eq:coords} maps the open
nonempty wedge $\varphi(\lambda)<\kappa<0$, $\lambda>0$, to an open
nonempty region in the original parameter plane.
\end{proof}

\begin{remark}[Terminology and scope]
The discriminant curve here is a fold tangent to the boundary
$\kappa=0$ at $\lambda=0$. The theorem does not identify it with a
Whitney cusp. It classifies every critical point in a fixed small radius
interval for nearby parameters, not every critical point on the whole
unit disk. In particular it does not refute the manuscript's separate
integer-outer-order conjecture for $a\ge2$.
\end{remark}

\subsection{An explicit asymptotic two-turning-point family}
\begin{corollary}[Separated small-radius extrema]
\label{rbif:cor:twofamily}
There is an analytic order pair $(a(\varepsilon),b(\varepsilon))$ for
small $\varepsilon$ with
\begin{equation}
 K=6\rbrr\varepsilon^2,\qquad
 Q=-\frac92\rbrr\varepsilon.
 \label{rbif:eq:epsfamily}
\end{equation}
For every sufficiently small $\varepsilon>0$, its normalized-radius
curve has exactly two critical points in the local interval of
Theorem~\ref{rbif:thm:unfolding}, with
\begin{equation}
 \rho_{\min}=\sqrt\varepsilon\,(1+O(\varepsilon)),\qquad
 \rho_{\max}=\sqrt{2\varepsilon}\,(1+O(\varepsilon)).
 \label{rbif:eq:epsroots}
\end{equation}
\end{corollary}
\begin{proof}
Analytic inversion of $(a,b)\mapsto(K,Q)$ defines the family.
Substitute $t=\varepsilon s$ into \eqref{rbif:eq:U}. Uniformly for
bounded $s$,
\[
 \varepsilon^{-2}U
 =3\rbrr(s-1)(s-2)+O(\varepsilon).
\]
At each of $s=1,2$ the limiting root is simple. The implicit-function
theorem gives roots $s=1+O(\varepsilon)$ and
$s=2+O(\varepsilon)$, with the stated types. The complete local
classification excludes additional local critical points.
\end{proof}

\subsection{The fourth-root onset law}
\begin{corollary}[Fixed-inner-order onset at the transition]
\label{rbif:cor:fourthroot}
Fix $b=\rbpb$ and let $a\uparrow\rbpa$ from below. There is exactly
one small-radius critical point, a strict maximum, and
\begin{equation}
 \rho_{\max}(a,\rbpb)
 =\left(\frac{K_a(\rbpa,\rbpb)}{3\rbrr}\right)^{1/4}
       (\rbpa-a)^{1/4}
       +O\bigl((\rbpa-a)^{3/4}\bigr).
 \label{rbif:eq:fourthroot}
\end{equation}
For $a>\rbpa$ sufficiently close there is no small-radius critical
point. The ratio inside the fourth root is positive.
\end{corollary}
\begin{proof}
Let $h=\rbpa-a>0$. Then
$\kappa=-K_a(\rbpa,\rbpb)h+O(h^2)>0$ and $\lambda=O(h)$.
Set $t=\sqrt h\,s$ in \eqref{rbif:eq:U}; after division by $h$ the
limiting equation is
$-K_a(\rbpa,\rbpb)+3\rbrr s^2=0$.
Its positive root is simple. Applying the implicit-function theorem in
$\sqrt h$ and then taking a positive square root of $t$ proves
\eqref{rbif:eq:fourthroot}. Uniqueness and maximality follow from
Theorem~\ref{rbif:thm:unfolding}.
For $a>\rbpa$, $\kappa<0$ and $\lambda=O(|a-\rbpa|)$.
If $\lambda\le0$ there is no root by the first row of the table.
If $\lambda>0$, then
$\varphi(\lambda)=O((a-\rbpa)^2)$, whereas $\kappa$ is negative
of first order. Thus $\kappa<\varphi(\lambda)$ for sufficiently small
positive $a-\rbpa$, and there is again no root.
\end{proof}

The previously proved maxima and minima at nonzero $q(b)$ have a
square-root radius law. The exponent $1/4$ is a genuine change of local
scaling, forced by the vanishing quartic coefficient and the nonzero
sixth-order coefficient; it is not a fit to a numerical curve.

\section{A certified rational witness}
\label{rbif:sec:witness}
The preceding existence theorem has unspecified neighborhood sizes.
The following result gives a concrete order pair and a fully explicit
radius interval. Each terminating decimal in the statement denotes the
exact rational number, not an approximation to an implicitly defined pair.

\begin{theorem}[A rational curve with at least two turning points]
\label{rbif:thm:witness}
Take
\begin{align*}
 a_0&=1.0014300032325009280102804945924913210012973562600080008779598088,\\
 b_0&=0.99749410117972379489824653149307623117370381382325680750587755176.
\end{align*}
The branch $\eta_{a_0,b_0}(\rho)$ has a strict local minimum in
\[
 \frac1{\sqrt{20000}}<\rho<\sqrt{\frac3{20000}}
\]
and a strict local maximum at a larger radius in
\[
 \sqrt{\frac3{20000}}<\rho<\frac1{\sqrt{4000}}.
\]
In particular, at least two positive-radius turning points occur below
$\rho=1/\sqrt{4000}$.
\end{theorem}
\begin{proof}
Use \eqref{rbif:eq:Psi} with $N=16$, so only
$p_2,\ldots,p_{35}$ are required. The exact interval replay constructs
an $\eta$ strip of half-width $10^{-8}$ centered at a rational enclosure
midpoint of $\mu(a_0,b_0)$. Throughout
$0\le t\le1/4000$, the left boundary has $\Psi<0$, the right boundary
has $\Psi>0$, and
\begin{equation}
 0.99900<\Psi_y<0.99906
 \label{rbif:eq:witnessstrip}
\end{equation}
throughout the strip. These inequalities use the infinite-series tails
\eqref{rbif:eq:tail}--\eqref{rbif:eq:tailt}, together with the Cauchy
bound for the $y$ derivative. There is therefore a unique analytic root
$y=\Phi(t)$ on this entire explicit interval. It is the branch from
Theorem~\ref{rbif:thm:generating}.

At each of three rational values of $t$, the replay encloses the root in
an interval of half-width $10^{-42}$ and uses
$\Phi'(t)=-\Psi_t/\Psi_y$. It certifies the following inequalities:
\begin{center}
\begin{tabular}{@{}cc@{}}
\toprule
$t$ & Exact enclosure for $\Phi'(t)$\\
\midrule
$1/20000$ & $-7.59\cdot10^{-19}<\Phi'(t)<-7.56\cdot10^{-19}$\\[1mm]
$3/20000$ & $ 2.52\cdot10^{-19}<\Phi'(t)< 2.54\cdot10^{-19}$\\[1mm]
$1/4000$ & $-7.59\cdot10^{-19}<\Phi'(t)<-7.56\cdot10^{-19}$\\
\bottomrule
\end{tabular}
\end{center}
The verifier is \path{verify_witness.py}; exact endpoints, all tail
bounds, the root strip, and the rational order pair are retained in
\path{witness_verified.json}. Since $\partial_\rho\eta=2\rho\Phi'(t)$,
the signs in the physical radius are negative, positive, negative.
Analyticity and the nonzero endpoint derivatives imply a strict local
minimum between the first two samples and a strict local maximum between
the last two. This proves the assertion for the exact infinite
polylogarithm.
\end{proof}

\begin{remark}[What this certificate does not count]
The witness proves at least two extrema in the displayed interval. It
does not claim that this particular rational pair has exactly two
critical points on $0<\rho\le1$. The local classification proves an exact
count in its own uniform small neighborhood; the explicit witness is
independently certified and does not assume that its parameters lie in
an unspecified neighborhood from an asymptotic theorem.
\end{remark}

The witness was found by taking $\varepsilon=10^{-4}$ in
\eqref{rbif:eq:epsfamily}, solving numerically for an order pair, and
rounding it to the displayed rational values. This history is not used
in the proof: the replay begins with those rational values and proves
the three signs afresh. The tiny derivative sizes also explain why an
ordinary double-precision plot would not be reliable evidence.

\section{Rigidity of the constant curve and an order-nine defect}
\label{rbif:sec:rigidity}
\begin{theorem}[Constant-curve classification]
\label{rbif:thm:rigidity}
For $a\ge1$ and $b>0$, the normalized-radius branch is constant on a
neighborhood of $\rho=0$ if and only if $(a,b)=(1,1)$.
In that case
\begin{equation}
 F_{1,1}(z)=\frac12\log^2(1-z),\qquad
 \eta_{1,1}(\rho)=\frac12\quad(0<\rho\le1).
 \label{rbif:eq:classical}
\end{equation}
Equivalently, within this order range, a M\"obius symmetry
$F_{a,b}(T_c(z))=F_{a,b}(z)$ can occur only for $(a,b,c)=(1,1,1/2)$.
\end{theorem}
\begin{proof}
A constant curve has $K=Q=R=0$. The quadratic classification leaves
only $(a,b)=(1,1)$ or a pair $(A(b),b)$ with $0<b<1$.
In the latter case the quartic classification forces $b=\rbpb$.
But Theorem~\ref{rbif:thm:nondeg} gives $R(\rbpa,\rbpb)<0$, a
contradiction.

For $(1,1)$, termwise differentiation of \eqref{rbif:eq:def} gives
\[
 F_{1,1}'(z)=-\frac{\log(1-z)}{1-z}.
\]
The zero initial value proves the first identity in
\eqref{rbif:eq:classical}. On the upper disk, the imaginary part of the
logarithm is nonzero. Thus $\rbIm\log^2(1-z)=0$ exactly when
$|1-z|=1$, which is $\cos\theta=\rho/2$. This proves the second
identity, including the boundary by continuation at this zero.
The symmetry characterization follows from
Corollary~\ref{rbif:cor:flatstructure}.
\end{proof}

\begin{corollary}[A quantified obstruction to the next flat curve]
\label{rbif:cor:order9}
At $(\rbpa,\rbpb)$, set $c=\rbmu$ and $w=z/(1-cz)$. Then
\begin{equation}
 F_{\rbpa,\rbpb}(z)-F_{\rbpa,\rbpb}(T_c(z))
 =-4\,2^{-\rbpa}\rbrr\,w^9+O(w^{11}).
 \label{rbif:eq:order9}
\end{equation}
The leading coefficient is strictly positive. The transformed function
is even through degree eight but fails to be even at degree nine.
\end{corollary}
\begin{proof}
At $c=\mu$, $P_0(c)=0$. Since $K=Q=0$,
\eqref{rbif:eq:K}--\eqref{rbif:eq:Q} also give $P_1(c)=P_2(c)=0$.
The transport identity then gives $d_3=d_5=d_7=0$ and
$d_9=P_3(c)=-2p_2\rbrr$.
Taking the odd part of $G_c$ proves \eqref{rbif:eq:order9}, and
$p_2=2^{-\rbpa}>0$, $\rbrr<0$ determine the sign.
\end{proof}

This supplies a structural explanation for the nearly flat transition:
it satisfies the first three odd symmetry conditions but not the fourth.
The functional relation \eqref{rbif:eq:mobiusPsi} makes the explanation
exact, rather than merely an analogy between a Taylor expansion and a
geometric picture.

\section{Audit, proposed manuscript updates, and limitations}
\label{rbif:sec:audit}
The formulas for $K$ and $Q$ in the audited local-radius and turning
sections agree with the independent derivations here. We found no
incorrect formula in those inspected portions. The appropriate
integration is therefore a theorem-status update, not a claim that the
previously cautious statements were errors.

The closing paragraph of \path{05-real-turning.tex} can now be updated
in three precise ways. The number of zeros of $q$ in $(0,1)$ is one;
the displayed decimal location can be replaced by the certified intervals
\eqref{rbif:eq:bbox}--\eqref{rbif:eq:abox}; and the higher-order behavior
can cite the negative sixth-order coefficient, the two-turning-point
region, and the fourth-root onset law. The earlier separate maximum and
minimum theorems should remain credited as predecessors and useful
special cases of Corollary~\ref{rbif:cor:localfamily}.

The source's warning that a nonzero local coefficient does not determine
whole-disk monotonicity remains important. Our two-turning-point theorem
strengthens that warning. It does not prove the full-radius conjecture
for integer outer orders $a\ge2$, classify all fractional curves on the
unit disk, or extend the constant-curve theorem to $0<a<1$.
Nothing here proves the manuscript's separate $S_6$ or $S_8$ arithmetic
candidates, numerical period independence, or a minimal-depth claim.

All finite sign decisions in the main certificates use only Python
integers and rational input values. The optional symbolic and numerical
regressions are distinguished from the proof. There is no Lean
formalization, and successful replay does not replace independent review
of the analytic reductions and interval implementation. The trust base
is explicit: the written analytic arguments, Python integer arithmetic,
the interval code, and the supplied rational boxes.

The package is designed for placement under
\path{Analysis/Polylogarithms/docs/reports/radial-boundary-bifurcation/}.
It contains a namespaced manuscript excerpt, bibliography additions,
a narrow proposed status patch, provenance, and replay instructions.
No repository files were modified as part of preparing this article.

\section{Further research questions}
\label{rbif:sec:future}
\begin{enumerate}[label=\textbf{\arabic*.},leftmargin=*]
\item \textbf{Whole-disk turning-point counts.}
Determine whether the explicit rational witness has exactly two critical
points on $0<\rho\le1$. More generally, classify the continuation of the
local two-turning-point region up to the unit circle. A local fold does
not by itself exclude further critical points at larger radii.
\item \textbf{The integer-outer-order conjecture.}
Prove or disprove strict decrease of $\eta_{a,b}$ on the whole interval
for $a\ge2$, $b>0$, and the proposed $a=1$ alternatives. The present
counterexamples all have fractional outer order strictly between one
and two, so the integer conjecture survives unchanged.
\item \textbf{Constant curves below outer order one.}
Extend Theorem~\ref{rbif:thm:rigidity} to $0<a<1$. The M\"obius normal
form reduces the question to an infinite, triangular family of exact
coefficient identities. The first three obstructions already suffice
for $a\ge1$, but a corresponding classification of $K=0$ below one is
not supplied here.
\item \textbf{An analytic proof of quartic uniqueness.}
Replace the $2304$-cell certificate by a direct inequality for
$q'(b)$ on $[0,9/10]$ and for $q''(b)$ on $[9/10,1]$, or find a simpler
single transformed quantity with a fixed derivative sign. The present
certificate suggests strict margins but is not such a symbolic proof.
\item \textbf{Higher contact and multiple polylogarithms.}
For deeper real-order polylogarithms, derive the analogue of
\eqref{rbif:eq:mobiusPsi} and investigate how many independent odd
symmetry conditions can vanish. Determine whether higher-dimensional
order families produce radial onset exponents beyond $1/4$.
\item \textbf{Sharp quantitative neighborhoods.}
Give explicit $t_0$ and parameter radii for every row of
Theorem~\ref{rbif:thm:unfolding}. An interval Taylor model adapted to
$(K,Q)$ coordinates should be more efficient than a raw rectangle in
$(a,b)$ because the coefficient map is poorly conditioned.
\item \textbf{Formal and cross-system verification.}
Formalize the interval arithmetic and the centered implicit-derivative
bounds in a proof assistant, or independently replay the analytic
inequalities with a separately implemented interval system. This would
reduce implementation risk without conflating finite checking with
formalization of the entire analytic argument.
\end{enumerate}

\appendix
\section{The finite verification argument in reproducible form}
\label{rbif:app:verification}
\subsection{Outward arithmetic and elementary functions}
Every interval is represented by integer endpoints $[L,U]$ denoting
$[L2^{-256},U2^{-256}]$. Addition and negation are exact on this grid.
For multiplication, the four integer endpoint products are computed;
the least is divided by $2^{256}$ with floor rounding and the greatest
with ceiling rounding. A positive reciprocal interval is obtained from
$2^{512}/U$ and $2^{512}/L$ with outward rounding. Negative reciprocals
are reduced by sign reversal. Division by an interval containing zero
is rejected. The resulting rules are inclusion-preserving by elementary
monotonicity, so an induction over an expression proves that its
computed interval contains its exact range.

For $2\le n\le9$, use
\begin{equation}
 \log n=2\sum_{j=0}^{N-1}\frac{z^{2j+1}}{2j+1}+E_N,\qquad
 0<E_N<\frac{2z^{2N+1}}{(2N+1)(1-z^2)},\quad
 z=\frac{n-1}{n+1},\quad N=512.
 \label{rbif:eq:logtail}
\end{equation}
The witness also needs $n\le35$. For $n>9$, write $n=2^km$ with
$1\le m<2$, use the same formula for $\log m$ with $N=128$, and add
$k\log2$. Here $z=(n-2^k)/(n+2^k)<1/3$.
For $0\le x\le8$, set $y=x/16$ and use
\begin{equation}
 e^y=\sum_{j=0}^{64}\frac{y^j}{j!}+E,\qquad
 0\le E\le\frac{y^{65}}{65!}\frac1{1-y/66}.
 \label{rbif:eq:exptail}
\end{equation}
Four interval squarings enclose $e^x$, and reciprocation encloses
$e^{-x}$. All exponential arguments in the certificates lie in the
allowed interval. The logarithm and exponential tail bounds are positive
geometric majorants; no floating-point library supplies a proof bound.

\subsection{Third-order derivatives along the threshold}
The implementation stores the ten Taylor coefficients with multi-indices
$(i,j)$ satisfying $i+j\le3$. Entries are divided derivatives
$\partial_a^i\partial_b^j/(i!j!)$, so multiplication is ordinary truncated
convolution. Inversion is obtained recursively from the identity
$XX^{-1}=1$. The derivatives of each normalized coefficient can also
be checked directly from
\[
 r_n=(2/n)^a\sum_{m<n}m^{-b},\qquad
 \partial_a^i\partial_b^j r_n
 =(-\log(n/2))^i(2/n)^a\sum_{m<n}(-\log m)^j m^{-b}.
\]
This gives a short independent specification of every jet input.

Let $v=A'$, $w=A''$, as in \eqref{rbif:eq:implicitdiff}, and define
\begin{equation}
 \mathcal C(X)=X_{bbb}+3X_{abb}v+3X_{aab}v^2+X_{aaa}v^3
                   +3(X_{ab}+X_{aa}v)w.
 \label{rbif:eq:Cderiv}
\end{equation}
Then
\[
 A'''=-\frac{\mathcal C(K)}{K_a},\qquad
 q'''=\mathcal C(Q_0)+(Q_0)_aA'''.
\]
These are obtained by differentiating $K(A(b),b)=0$ and
$q(b)=Q_0(A(b),b)$ three times. They are used only after a strict
negative enclosure for $K_a$ has been established.

\subsection{Why the mesh is exhaustive}
The input file \path{global_mesh.json} is a list of rational proposals.
The verifier first checks that the first endpoint is zero, the last is
one, adjacent endpoints agree exactly, all cell lengths are positive,
and each listed midpoint is the exact arithmetic midpoint. It checks
that every cell lies in the correct side of $9/10$. It then verifies both
signs for every root bracket, checks the barrier derivatives on every
cell rectangle, and checks that the midpoint bracket lies within its
rectangle. Thus missing intervals, unverified root guesses, and
unjustified graph interpolation are not accepted.

The centered derivative step is essential. A natural interval evaluation
of $q'$ on a moderately wide rectangle loses too much correlation between
$a$ and $b$. Instead the exact midpoint value is combined with a bound on
one higher total derivative along the implicitly defined graph. The
mean-value theorem supplies the rigorous error. The same method one order
higher proves concavity near $b=1$.

\subsection{Replay and artifacts}
From the package root, the decisive checks are
\begin{verbatim}
python code/verify_global.py
python code/verify_local.py
python code/verify_witness.py
\end{verbatim}
They require only the Python standard library. They must not be run with
\texttt{-O}; optimized execution is explicitly rejected because the
verifiers use assertions for proof obligations. The optional check
\begin{verbatim}
python code/verify_symbolic.py
\end{verbatim}
requires SymPy and verifies the coefficient identities at finite orders,
the sixth-order formula, and the Bernstein expansion. It is a regression
check, not an extrapolation argument. The universal identities are proved
in the text. The optional generator \path{make_mesh.py} requires mpmath;
regenerating boxes is unnecessary to verify the delivered certificate.

\begin{center}
\begin{tabular}{@{}p{0.48\textwidth}p{0.43\textwidth}@{}}
\toprule
Artifact & Role\\
\midrule
\path{global_mesh.json} & Rational graph-cover proposals.\\[1mm]
\path{global_verified.json} & All 2304 strict cell margins and coverage receipt.\\[1mm]
\path{local_verified.json} & Contraction, location, Jacobian and sixth-order signs.\\[1mm]
\path{witness_verified.json} & Rational orders, analytic tails, root strip and three signs.\\[1mm]
\path{symbolic_verified.json} & Optional exact symbolic regression receipt.\\[1mm]
\path{integration/} & Namespaced excerpt, bibliography and proposed status replacement.\\[1mm]
\path{provenance/} & Audited repository paths, blob identifiers and build environment.\\
\bottomrule
\end{tabular}
\end{center}
The final archive also includes editable LaTeX, a compiled PDF, the source
code, and a checksum manifest. Receipts record successful replay for this
package; they do not automatically validate subsequently edited files.

\section{Proof dependencies and what has actually been established}
\label{rbif:app:dependencies}
The logic is deliberately separated into four layers.
First, the generating equation, reversion formula, M\"obius transport,
and local bifurcation arguments are ordinary analytic or formal proofs.
Second, the known quadratic classification supplies the globally defined
threshold $A(b)$; its proof is reproduced and credited.
Third, finitely many exact inequalities prove derivative separation,
local contraction, and an explicit witness. Fourth, the analytic
arguments transfer those inequalities to uniqueness, local extrema,
scaling laws, and rigidity for the infinite function.

The new global statement is global in the inner-order parameter
$0<b<1$ \emph{along the quadratic threshold}. The new exact turning-point
count is local in disk radius and in the two order parameters.
The concrete rational witness supplies an explicit small radius interval,
but only an at-least-two count. These three scopes are different and
should not be merged during manuscript integration.

The central resolved question can be summarized without any numerical
notation: along the unique quadratic threshold in the strip
$1<a<2$, $0<b<1$, the quartic coefficient changes sign exactly once;
at that change the sixth-order coefficient is negative. This finite
coefficient fact forces an open two-extremum region and rules out a second
constant normalized-radius curve for every $a\ge1$, $b>0$.

\begin{thebibliography}{9}
\bibitem{rbif:local}
ProveIt Contributors,
\emph{A proved local part of the normalized-radius conjecture},
\path{Analysis/Polylogarithms/docs/manuscript/chapters/05-real-local-radius.tex},
Git blob \texttt{30e3e31944be5bc0d03a8485386dd1608c075507},
audited October 10, 2026.
Repository: \url{https://github.com/VladimirReshetnikov/ProveIt}.

\bibitem{rbif:turning}
ProveIt Contributors,
\emph{A proved failure of fractional global monotonicity},
\path{Analysis/Polylogarithms/docs/manuscript/chapters/05-real-turning.tex},
Git blob \texttt{49f0336656b6aa893e5ba6e904fb561c3a941cfb},
audited October 10, 2026, at repository commit
\texttt{0707425e155707e53d2892e30c72e54d06c00e5f}.

\bibitem{rbif:cartesian}
ProveIt Contributors,
\emph{Cartesian motion of the complete real-order zero arc},
\path{Analysis/Polylogarithms/docs/manuscript/chapters/05-real-cartesian-motion.tex},
Git blob \texttt{0deb9bb54adaef8783d4e0e7ce43d426643450a0},
audited October 10, 2026.

\bibitem{rbif:gessel}
Ira M. Gessel,
\emph{Lagrange inversion},
Journal of Combinatorial Theory, Series A \textbf{144} (2016), 212--249.
\url{https://doi.org/10.1016/j.jcta.2016.06.018}.

\bibitem{rbif:abl}
Jakob Ablinger, Johannes Bl\"umlein, and Carsten Schneider,
\emph{Harmonic sums and polylogarithms generated by cyclotomic polynomials},
Journal of Mathematical Physics \textbf{52} (2011), 102301.
\url{https://arxiv.org/abs/1105.6063}.
\end{thebibliography}

\end{document}
